% Options for packages loaded elsewhere
\PassOptionsToPackage{unicode}{hyperref}
\PassOptionsToPackage{hyphens}{url}
%
\documentclass[
]{article}
\usepackage{amsmath,amssymb}
\usepackage{lmodern}
\usepackage{iftex}
\ifPDFTeX
  \usepackage[T1]{fontenc}
  \usepackage[utf8]{inputenc}
  \usepackage{textcomp} % provide euro and other symbols
\else % if luatex or xetex
  \usepackage{unicode-math}
  \defaultfontfeatures{Scale=MatchLowercase}
  \defaultfontfeatures[\rmfamily]{Ligatures=TeX,Scale=1}
\fi
% Use upquote if available, for straight quotes in verbatim environments
\IfFileExists{upquote.sty}{\usepackage{upquote}}{}
\IfFileExists{microtype.sty}{% use microtype if available
  \usepackage[]{microtype}
  \UseMicrotypeSet[protrusion]{basicmath} % disable protrusion for tt fonts
}{}
\makeatletter
\@ifundefined{KOMAClassName}{% if non-KOMA class
  \IfFileExists{parskip.sty}{%
    \usepackage{parskip}
  }{% else
    \setlength{\parindent}{0pt}
    \setlength{\parskip}{6pt plus 2pt minus 1pt}}
}{% if KOMA class
  \KOMAoptions{parskip=half}}
\makeatother
\usepackage{xcolor}
\IfFileExists{xurl.sty}{\usepackage{xurl}}{} % add URL line breaks if available
\IfFileExists{bookmark.sty}{\usepackage{bookmark}}{\usepackage{hyperref}}
\hypersetup{
  hidelinks,
  pdfcreator={LaTeX via pandoc}}
\urlstyle{same} % disable monospaced font for URLs
\usepackage{longtable,booktabs,array}
\usepackage{calc} % for calculating minipage widths
% Correct order of tables after \paragraph or \subparagraph
\usepackage{etoolbox}
\makeatletter
\patchcmd\longtable{\par}{\if@noskipsec\mbox{}\fi\par}{}{}
\makeatother
% Allow footnotes in longtable head/foot
\IfFileExists{footnotehyper.sty}{\usepackage{footnotehyper}}{\usepackage{footnote}}
\makesavenoteenv{longtable}
\usepackage{graphicx}
\makeatletter
\def\maxwidth{\ifdim\Gin@nat@width>\linewidth\linewidth\else\Gin@nat@width\fi}
\def\maxheight{\ifdim\Gin@nat@height>\textheight\textheight\else\Gin@nat@height\fi}
\makeatother
% Scale images if necessary, so that they will not overflow the page
% margins by default, and it is still possible to overwrite the defaults
% using explicit options in \includegraphics[width, height, ...]{}
\setkeys{Gin}{width=\maxwidth,height=\maxheight,keepaspectratio}
% Set default figure placement to htbp
\makeatletter
\def\fps@figure{htbp}
\makeatother
\setlength{\emergencystretch}{3em} % prevent overfull lines
\providecommand{\tightlist}{%
  \setlength{\itemsep}{0pt}\setlength{\parskip}{0pt}}
\setcounter{secnumdepth}{-\maxdimen} % remove section numbering
\ifLuaTeX
  \usepackage{selnolig}  % disable illegal ligatures
\fi

\author{}
\date{}

\begin{document}

\hypertarget{ux95eeux9898ux4e09ux89c6ux89c9ux8ba4ux77e5ux5b8fux89c2ux8111ux7535ux6a21ux578bux4e0eux5b9eux6d4bux9a8cux8bc1}{%
\section{5.3
问题三：视觉认知宏观脑电模型与实测验证}\label{ux95eeux9898ux4e09ux89c6ux89c9ux8ba4ux77e5ux5b8fux89c2ux8111ux7535ux6a21ux578bux4e0eux5b9eux6d4bux9a8cux8bc1}}

\hypertarget{ux95eeux9898ux5206ux6790ux4e0eux5efaux6a21ux8defux7ebf}{%
\subsection{5.3.1
问题分析与建模路线}\label{ux95eeux9898ux5206ux6790ux4e0eux5efaux6a21ux8defux7ebf}}

问题三要求在问题二脑电形成机制的基础上建立认知宏观模型，并用题目提供的
EEG 验证模型；同时需考虑应答可能错误或未及时发生{[}1{]}。题目数据记录了
F3、Fz、F4 三个原始脑电通道、视觉提示通道 VisCue、应答通道 Action/TgtAct
及时间戳。三路额区 EEG 是有限的头皮观测，不能据此唯一反演
LGN、海马或前额叶的脑源活动。因此，本文把问题二的视觉通路和源---场映射作为认知模型的机制约束，但第三问的数值分析直接使用题目提供的原始
EEG，不把问题二生成的模拟信号当作第三问的观测数据。

据此将第三问分为三个相互区分的分析环节：首先从原始 EEG
建立试次级观测特征；其次以视觉前馈、记忆匹配和认知控制为解释框架，构造与特征一一对应的
V/H/P 宏观代理状态；最后分别检验 EEG 对提示方向的跨记录预测能力、V/H/P
代理的消融表现，以及应答侧是否能由 EEG
状态补充解释。由于当前数据没有经独立核验的目标出现时刻、反应截止时间和逐试次正确答案，本文不把应答选择分类称为正确率，也不在不合理的反应时上拟合漂移扩散模型（DDM）。

整体计算链为：

\begin{verbatim}
原始 F3/Fz/F4 连续 EEG
       ↓ 连续滤波、事件切段与质量标记
试次级 ERP / 频带功率特征
       ↓ 固定特征定义
V/H/P 可观测功能代理与描述性路径
       ↓ 按完整记录留出
提示方向验证、代理消融与应答侧补充分析
\end{verbatim}

\hypertarget{ux6570ux636eux5904ux7406ux4e0eux8bd5ux6b21ux7ea7ux7279ux5f81}{%
\subsection{5.3.2
数据处理与试次级特征}\label{ux6570ux636eux5904ux7406ux4e0eux8bd5ux6b21ux7ea7ux7279ux5f81}}

四份 MAT 记录的采样率均为 256 Hz，每份记录含 100 个 VisCue 事件，共 400
个候选试次。所有 EEG 特征仅由原始 F3、Fz、F4 计算；数据集自带的
F3Decon、FzDecon、F4Decon 不进入特征矩阵。VisCue 提供左、右提示标签
\(u_i\in\{-1,+1\}\)。Action/TgtAct 与时间戳只用于应答元数据，不进入 EEG
特征。程序保留原始应答值，并统一生成标准码：原码 \texttt{-1/-2} 映射为
\texttt{response\_code=-2}（左），原码 \texttt{+1/+2} 映射为
\texttt{response\_code=+2}（右）；对检出应答边沿的试次，令应答侧
\(a_i=\text{response\_code}/2\)，得到左 \(-1\)、右 \(+1\)。通道 9 只检测
\(0\to\text{非零}\)
的起始沿，每个起始沿计作一次应答，持续的非零段不重复计数。四份记录共检出
400 个应答边沿，边沿原码为 \texttt{-1} 的有 196 个、\texttt{+1} 的有 204
个；连续信号中出现的少量 \texttt{±2}
采样点不是额外的应答边沿。若试次没有检出边沿，则 \texttt{response\_code}
记为 0，原始 \texttt{response\_raw} 留空，应答侧保持缺失。

原始连续 EEG 分为 ERP 与时频两条支路。ERP 支路使用四阶零相位 Butterworth
带通滤波（0.5--30 Hz），时频支路使用四阶零相位 Butterworth
带通滤波（1--80
Hz）；滤波在整段连续记录上完成，再按事件切段。提示锁定窗为
({[}-0.10,0.50)) s；目标锁定窗为 ({[}-0.10,0.80)) s。题目说明提示持续约
0.2 s、随后等待约 2 s，因此程序暂用 (t\_\{target\}=t\_\{cue\}+2.2) s
作为目标时刻，并另检查 2.0--2.4 s
的时间偏移。该目标时刻来自实验排程推定，不是独立检测到的目标事件标记。

ERP 特征以提示前 \([-0.10,0)\) s 的均值为基线，对 250--500 ms
区间作基线校正。记第 \(i\) 个试次、通道 \(c\) 的信号为
\(x_{i,c}(t)\)，则额区 ERP 候选均值为

\[
E_{i,c}=\frac{1}{|W_E|}\sum_{t\in W_E}x_{i,c}(t)-\frac{1}{|W_B|}\sum_{t\in W_B}x_{i,c}(t),
\quad W_B=[-0.10,0),\quad W_E=[0.25,0.50).
\]

频带功率由事件后窗内的 Welch 功率谱密度积分得到，并取自然对数：

\[
L_{i,c,b}=\log\left(\int_{f_{b,\min}}^{f_{b,\max}}\widehat S_{i,c}(f)\,df\right),
\]

其中 \(b\) 分别为 theta（4--8 Hz）、alpha（8--13 Hz）、beta（13--30
Hz）和 gamma（30--80 Hz）。提示锁定功率窗为 0--500
ms，目标锁定功率窗为目标事件后 0--800
ms。提示窗较短，低频功率估计的频率分辨率有限；gamma 仅作为探索性特征。

用于提示方向分类的 13 维 EEG 特征向量为

\[
\mathbf f_i=\big[E_{i,F3},E_{i,Fz},E_{i,F4},
L_{i,F3,\theta},L_{i,Fz,\theta},L_{i,F4,\theta},
L_{i,F3,\alpha},L_{i,Fz,\alpha},L_{i,F4,\alpha},
L_{i,F3,\beta},L_{i,Fz,\beta},L_{i,F4,\beta},
A_{i,\alpha}\big]^T,
\]

其中 \(A_{i,\alpha}=L_{i,F4,\alpha}-L_{i,F3,\alpha}\)。13 维特征由 3 个
ERP 均值、9 个三通道频带功率及 1 个额区 alpha
对数功率差组成。程序另外提取 ERP
峰值、峰潜伏期、额区不对称和应答终点窗口特征；这些变量不全部纳入主提示分类器，避免无依据地扩大特征空间。

应答终点支路以通道 9 检出的应答标记时刻 (t\_\{act\})
为终点，提取从提示开始至 (t\_\{act\}-0.10) s 的累计信号，并计算末端 500
ms
汇总特征。该支路用于探索题目指定的应答前信号，不并入提示方向分类器；在应答标记语义尚未核实前，也不将它解释为经行为反应时校准的认知终点。

质量控制只标记，不插值或改写信号：若试次窗含非有限值、通道平线，或任一绝对值达到原始数据截断阈值
999.5，则该窗判为无效。没有独立眼电参考且仅有三路额区
EEG，故不宣称实施了 ICA 眼动去除。原始单位未由数据说明单独确认，ERP
数值以下均称为原始数据单位。

\begin{longtable}[]{@{}lrrl@{}}
\toprule
分析阶段 & 候选试次数 & 通过原始 EEG 质量控制 & 用途 \\
\midrule
\endhead
提示锁定 & 400 & 387 & 主提示方向验证及 V/H/P 代理建模 \\
目标锁定（偏移 2.2 s） & 400 & 396 & 次要阶段比较与时间敏感性分析 \\
应答标记前窗口 & 400 & 381 & 应答前信号描述，不作为主分类输入 \\
Q1 质量匹配敏感性样本 & --- & 297 & 检查前序质量筛选口径的影响；不作为
Q3 EEG 特征来源 \\
\bottomrule
\end{longtable}

图 5.3-1
展示目标锁定特征在左右提示下的分布。箱体重叠意味着单变量分布不能单独支持稳定的左右效应；该图用于描述，不用于筛选特征或调分类器。

\begin{figure}
\centering
\includegraphics{../output/features_by_cue_side.png}
\caption{图 5.3-1 左右提示下的目标锁定 EEG 特征分布}
\end{figure}

\hypertarget{vhp-ux8bd5ux6b21ux7ea7ux5b8fux89c2ux4ee3ux7406ux6a21ux578b}{%
\subsection{5.3.3 V/H/P
试次级宏观代理模型}\label{vhp-ux8bd5ux6b21ux7ea7ux5b8fux89c2ux4ee3ux7406ux6a21ux578b}}

题目提出的视觉前馈、记忆匹配和认知控制分别记为
(V\_i,H\_i,P\_i)。在三通道头皮 EEG 上拟合自由稀疏观测矩阵 (C)
会面临源混合与状态旋转不可辨识问题；当前程序因此不估计自由潜变量载荷，也不把代理变量称为
LGN、海马或 PFC 的直接活动。程序用三组可重复计算的 EEG
特征构造功能代理：

\[
V_i=\frac{1}{3}\sum_{c\in\{F3,Fz,F4\}}E_{i,c},\qquad
H_i=\frac{1}{3}\sum_c L_{i,c,\theta},\qquad
P_i=\frac{1}{3}\sum_c L_{i,c,\beta}.
\]

其中 (V\_i) 为额区 ERP 候选均值，(H\_i) 为额区 theta
对数功率均值，(P\_i) 为额区 beta 对数功率均值。gamma
功率只用于扩展消融。上述定义是对可观测特征的功能性命名，不能据此定位深部来源。

在该代理层上，以提示方向 \(u_i\) 描述视觉输入与 \(V_i\)
的关系，并以记录文件名中的 Task-1/Task-2 后缀构造中心化候选变量
\(T_i^*=\text{task code}-1.5\)，拟合以下描述性路径：

\[
\begin{aligned}
V_i &= \alpha_V+\beta_Vu_i+\varepsilon_{V,i},\\
H_i &= \alpha_H+\beta_{HV}V_i+\beta_{HT}T_i^*+\varepsilon_{H,i},\\
P_i &= \alpha_P+\beta_{PH}H_i+\beta_{PV}V_i+\beta_{PT}T_i^*+\varepsilon_{P,i}.
\end{aligned}
\]

系数由全部质量合格试次作 OLS
描述性估计，没有按独立被试层级计算的显著性检验；(T\_i\^{}*)
的正式项目语义亦未由实验元数据核实。因此路径系数用于检查特征间的样本内关联，不作因果路径或脑区强度解释。提示侧
(u\_i) 只出现在上述描述性 (V\_i) 方程中，不作为 EEG-only
提示分类器的输入。

\hypertarget{ux9a8cux8bc1ux65b9ux6848}{%
\subsection{5.3.4 验证方案}\label{ux9a8cux8bc1ux65b9ux6848}}

\textbf{提示方向验证。} 以 VisCue 的左/右标签为目标变量，对 13 维真实
EEG 特征使用固定 L2 Logistic 回归。按完整 MAT
记录留一：每次留出一份记录测试、其余三份训练；所有标准化参数仅由该折训练数据估计。分别报告提示锁定与目标锁定结果。主指标为每个留出记录的平衡准确率（BA）和
AUC，准确率与宏平均 F1 作为补充。

\[
BA_r=\frac{1}{2}\left(\frac{TP_{L,r}}{N_{L,r}}+\frac{TP_{R,r}}{N_{R,r}}\right),\qquad
\overline{BA}=\frac{1}{4}\sum_{r=1}^{4}BA_r.
\]

四份记录是当前可用的外层分组；由于文件与独立被试的映射未确认，本文称其为留一记录验证，不称为留一被试验证。另做记录内随机折诊断，但不以其作为泛化证据，因为训练与测试会共享同一记录。

\textbf{代理消融与重建。} 在同一留一记录划分下比较 V-only、V/H/P
代理组合、加入候选任务码及 gamma
的方案，并进行逐特征留出重建。重建误差越低表示留出特征在该代理组合下越容易预测；该结果是特征结构的描述性证据，不等于隐源模型得到识别。AIC/BIC
仅作同一数据口径下的辅助拟合量，不作为显著性检验。

\textbf{应答侧与 DDM 门槛。} 将标准 \texttt{response\_code} 除以 2
得到应答侧 \(a_i\in\{-1,+1\}\)，在留一记录验证下比较``提示方向 +
文件名任务码候选''基线与加入 EEG
代理后的模型。该分析预测的是记录到的应答侧，不代表判断正确与否。DDM
需要可信的目标时刻、反应时、截止时间及正确/错误/未及时标签；上述信息不齐全时不拟合
DDM，也不将缺少 marker 直接认定为漏答。

\hypertarget{ux7ed3ux679cux4e0eux5206ux6790}{%
\subsection{5.3.5 结果与分析}\label{ux7ed3ux679cux4e0eux5206ux6790}}

\hypertarget{ux5b9eux6d4b-eeg-ux7684ux5de6ux53f3ux63d0ux793aux89e3ux7801}{%
\subsubsection{（1）实测 EEG
的左右提示解码}\label{ux5b9eux6d4b-eeg-ux7684ux5de6ux53f3ux63d0ux793aux89e3ux7801}}

提示锁定窗有 387 个质量合格试次，留出记录的 BA 均值为 0.460、AUC 均值为
0.483；目标锁定窗有 396 个合格试次，BA 均值为 0.522、AUC 均值为
0.526。下表保留每个留出方向，避免均值掩盖记录差异。提示锁定分类整体接近机会水平；目标锁定结果略高，但主要由
VisualCogB\_Task-1 一折拉高，其他三折 BA 为
0.478--0.510，且目标时刻本身是排程推定值，故不能据此声称稳定的目标诱发左右解码。

\begin{longtable}[]{@{}llrrrrr@{}}
\toprule
阶段 & 留出记录 & 测试数 & Accuracy & Balanced Accuracy & AUC &
Macro-F1 \\
\midrule
\endhead
提示锁定 & VisualCogA\_Task-1 & 100 & 0.490 & 0.491 & 0.516 & 0.490 \\
提示锁定 & VisualCogA\_Task-2 & 96 & 0.448 & 0.417 & 0.443 & 0.392 \\
提示锁定 & VisualCogB\_Task-1 & 98 & 0.480 & 0.480 & 0.497 & 0.479 \\
提示锁定 & VisualCogB\_Task-2 & 93 & 0.452 & 0.454 & 0.478 & 0.435 \\
提示锁定 & 四记录等权均值 & 387 & 0.467 & 0.460 & 0.483 & 0.449 \\
目标锁定 & VisualCogA\_Task-1 & 100 & 0.490 & 0.489 & 0.494 & 0.489 \\
目标锁定 & VisualCogA\_Task-2 & 98 & 0.500 & 0.478 & 0.452 & 0.471 \\
目标锁定 & VisualCogB\_Task-1 & 98 & 0.612 & 0.612 & 0.618 & 0.612 \\
目标锁定 & VisualCogB\_Task-2 & 100 & 0.510 & 0.510 & 0.538 & 0.465 \\
目标锁定 & 四记录等权均值 & 396 & 0.528 & 0.522 & 0.526 & 0.509 \\
\bottomrule
\end{longtable}

\begin{figure}
\centering
\includegraphics{../output/eeg_heldout_record_confusion.png}
\caption{图 5.3-2 留一记录 EEG 提示方向预测的混淆矩阵}
\end{figure}

将提示时间偏移作为敏感性分析时，目标锁定 BA 在 cue 后 2.0--2.4 s 范围为
0.483--0.541。该范围反映未标记目标时刻带来的不确定性，不是通过搜索选出的最佳时间点。对同一
297 个 Q1 质量匹配试次重算后，提示锁定 BA 为 0.487，目标锁定 BA 为
0.510；与旧 Q1-clean 基线相比变化分别为 +0.027 和
+0.003，说明质量口径会影响小幅数值，但没有形成强的跨记录提示解码证据。

\hypertarget{vhp-ux4ee3ux7406ux8defux5f84ux4e0eux72b6ux6001ux5206ux5e03}{%
\subsubsection{（2）V/H/P
代理路径与状态分布}\label{vhp-ux4ee3ux7406ux8defux5f84ux4e0eux72b6ux6001ux5206ux5e03}}

目标锁定下左右提示的三类 V/H/P 描述性代理分布见图
5.3-3。图中以全体目标锁定试次计算的 z
分数仅用于可视化；预测模型中的标准化仍在每个训练折内单独估计。该图中左右分布高度重叠，不能作为显著左右状态差异的证据。

\begin{figure}
\centering
\includegraphics{../output/state_proxy_scores.png}
\caption{图 5.3-3 目标锁定 V/H/P 功能代理的左右提示分布}
\end{figure}

表 5.3-2 汇总全体质量合格试次上的路径回归系数。提示锁定的 \(u_i\to V_i\)
系数为 0.851，目标锁定为 -0.191；而 \(V_i\to H_i\) 系数分别为 0.0023 和
0.00013。由于目标阶段与提示阶段的信号尺度、时间锚点和特征分布不同，加上系数是未标准化的样本内
OLS
估计，不能直接比较其大小并解释为神经通路增益。其用途是给出当前代理特征之间的拟合关系，而不是验证一个已识别的潜变量因果模型。

\begin{longtable}[]{@{}lrrrrrr@{}}
\toprule
阶段 & (\beta\_V) & (\beta\_\{HV\}) & (\beta\_\{HT\}) & (\beta\_\{PH\})
& (\beta\_\{PV\}) & (\beta\_\{PT\}) \\
\midrule
\endhead
提示锁定 & 0.8513 & 0.00235 & 0.1341 & 0.1029 & 0.00055 & -0.2101 \\
目标锁定 & -0.1905 & 0.00013 & -0.0687 & 0.1297 & 0.00124 & -0.2367 \\
\bottomrule
\end{longtable}

注：(T) 为文件名 Task-1/Task-2 后缀的中心化候选码；系数为原始尺度 OLS
描述值，不附加试次级显著性解释。

\hypertarget{ux4ee3ux7406ux6d88ux878dux5e94ux7b54ux4fa7ux548cux63a2ux7d22ux6027ux5206ux7c7bux5668ux6bd4ux8f83}{%
\subsubsection{（3）代理消融、应答侧和探索性分类器比较}\label{ux4ee3ux7406ux6d88ux878dux5e94ux7b54ux4fa7ux548cux63a2ux7d22ux6027ux5206ux7c7bux5668ux6bd4ux8f83}}

在 raw-QC 样本上，V-only 与完整 V/H/P 加任务码候选的提示 BA 分别为 0.494
和 0.489；目标 BA 分别为 0.497 和 0.530。加入 gamma 后，提示和目标 BA
分别为 0.475 和 0.522。逐特征留出重建 RMSE 从 V-only 的 1.026 降到完整
V/H/P 加任务码候选的 0.954（提示阶段），从 1.015 降到
0.856（目标阶段）。这说明组合代理在目标阶段对部分特征的重建更好，但分类提升并不一致，且所有数值只来自四个记录折，不能把重建改进写成海马或前额叶源活动得到验证。

应答侧留一记录 BA 为：提示/任务码候选基线 0.755、EEG 代理单独
0.490、提示/任务码与 EEG 代理合并 0.747。当前代码没有显示 EEG
代理在基线之外带来增益。并且任务后缀与提示---应答关系的正式含义尚未核实，所以该分析只能称为应答侧选择预测，不是行为正确率。

另对 49 个预设分类器配置做过嵌套留一记录探索：raw-QC 提示阶段 BA 从固定
Logistic 基线的 0.460 升至 0.521（+6.05 个百分点）。但是 Q1
质量匹配提示样本上变化为 -0.02
个百分点，目标阶段没有相同提升；四份记录也已在多轮开发中被查看。因此该
0.521
仅作为模型搜索敏感性结果，不作为第三问的主要确认性证据，也不称为完整
V/H/P 模型准确率。

\begin{figure}
\centering
\includegraphics{figures/F11_scope_sensitivity.png}
\caption{图 5.3-4 分类器搜索结果的样本口径敏感性}
\end{figure}

\hypertarget{ux8ba8ux8bbaux4e0eux7ed3ux8bba}{%
\subsection{5.3.6 讨论与结论}\label{ux8ba8ux8bbaux4e0eux7ed3ux8bba}}

题目提出的视觉前馈---记忆匹配---认知控制链条，可以作为第三问的机制组织框架；但在当前三路额区
EEG 与四份记录下，V/H/P 只能由 ERP
和频带功率构造为可观测代理。留一记录提示方向分类整体接近机会水平，代理组合的分类增益不稳定，应答侧模型中
EEG
也没有在提示/任务码候选基线上带来提升。因此，当前数据支持的是``可以建立并审计试次级
EEG 特征与功能代理''，而不是``已经从头皮 EEG 识别出海马---PFC
因果状态链''。

应答行为的限制尤其重要。四份记录中均检测到通道 9
的非零事件，但这只能说明记录中有 marker，不能单独证明没有漏答。程序按
(cue+2.2) s 估计目标时刻后，所得反应时中位数仅 0.0148 s，399/400 小于
0.10 s，且该时刻没有独立目标 onset
marker、正式截止时间也未提供。因而当前没有足够可信的 RT 和删失标签，DDM
不拟合；正确/错误/未及时应答也不从通道 9 单独推断。

综上，本文按``机制继承、数值解耦''的路线完成了原始 EEG
的试次级特征提取、V/H/P
功能代理构造、记录留出预测和消融诊断。当前结果显示，三通道额区 EEG
对左右提示的跨记录可解码性有限，V/H/P
代理尚不能被解释为深部源定位或经识别的因果状态。若后续获得经核验的任务/被试映射、逐试次目标真值、目标实际
onset、反应截止时间及漏答定义，可再开展正确性分层和 DDM
验证；在这些信息到位前，应将结论限制在当前可观测特征与选择标签范围内。

\textbf{代码对应关系：} \texttt{01\_audit\_events.py}
审核事件与记录；\texttt{02\_extract\_trials.py}
建立试次表；\texttt{02b\_preprocess\_raw.py} 处理连续原始 EEG
并切段；\texttt{03\_extract\_features.py}
提取试次特征；\texttt{04\_cognitive\_state.py} 构造 V/H/P
代理与描述性路径；\texttt{05\_behavior\_model.py} 比较应答侧 Logistic
模型并执行 DDM 门槛判断；\texttt{06\_validate.py} 完成留一记录 EEG
验证；\texttt{07\_ablation.py} 完成代理消融与重建比较。

\hypertarget{ux53c2ux8003ux6587ux732e}{%
\subsection{参考文献}\label{ux53c2ux8003ux6587ux732e}}

{[}1{]} 2026 年中国研究生数学建模竞赛 C
题：《服务于脑机接口与精神性疾病诊断的脑电图计算模型》；题目附件与实验数据说明。

\end{document}
